% Diapositiva  42


% Kevin
%=============================================
% Puntos 11, 12, 13: Mejoras futuras,
% conclusión del problema inverso y cierre formal.
%=============================================

\begin{frame}{Conclusión Final del Proyecto}

    \begin{columns}[T]

        \column{0.48\textwidth}

        \begin{block}{Mejoras Futuras}
            \scriptsize
            \begin{itemize}
                \item Incorporar \textbf{modelos físicos más complejos} de propagación.
                \vspace{0.1cm}
                \item Considerar \textbf{medios no homogéneos} (terreno realista).
                \vspace{0.1cm}
                \item Aumentar la \textbf{cantidad de sensores} para mayor resolución espacial.
                \vspace{0.1cm}
                \item Extender la estimación a \textbf{parámetros adicionales} del modelo sísmico.
            \end{itemize}
        \end{block}

        \column{0.48\textwidth}

        \begin{exampleblock}{El Problema Inverso se Resolvió}
            \scriptsize
            El problema inverso pudo resolverse correctamente utilizando:
            \begin{itemize}
                \item \textbf{Funciones multivariables} $E(x,y,z,A_0)$
                \vspace{0.05cm}
                \item \textbf{Gauss-Newton} — 8 iteraciones, precisión $>99.7\%$
                \vspace{0.05cm}
                \item \textbf{Análisis visual} de la función de error
                \vspace{0.05cm}
                \item \textbf{Estimación simultánea} de las 4 incógnitas
            \end{itemize}
        \end{exampleblock}

    \end{columns}

    \vspace{0.3cm}

    \begin{block}{}
        \centering
        \normalsize
        Este proyecto permitió integrar conceptos de \textbf{cálculo multivariado},
        \textbf{geometría espacial} y \textbf{optimización} en un problema aplicado
        de \textbf{localización sísmica}.
    \end{block}

    \vspace{0.2cm}
    \begin{center}
        \footnotesize
        \textbf{\textcolor{blue}{¡Gracias por su atención!}}
    \end{center}

\end{frame}